% Plan: development/outline.md, section 6, item C. Proposition prop:pts-exact and
% Remark rem:lukvle10-seeds moved here from Section 6 (development/moves/m-points-audit.md,
% Moves 1-3); the waterno2 p4 violation from development/moves/m-results.md, Move 3.
% Round-2 revision (adjudication G6-01, G6-10): "the authors' code" -> "the first
% code"; implementation details that the certificate box or the family sections of
% S1 state are not repeated; test records to development/moves/r2-G6.md.
\section{Exactly feasible points}\label{app:points}

All statements below concern the stored OSIL model under \reading{b}; the claim register (\cref{app:repro-register}) names the checkers of each point.

\subsection{Proof of \texorpdfstring{\cref{thm:pts-krawczyk}}{the Krawczyk theorem}}\label{app:points-krawczyk}

\begin{proof}[Proof of \cref{thm:pts-krawczyk}]
Let $w_1,\dots,w_m>0$ be the widths of $X$, and write $|\mathbf{B}_{ij}|=\max\{|b|:b\in\mathbf{B}_{ij}\}$.

\emph{Step 1 (mean-value form).}
Let $u,z\in X$.
The box $X$ is convex, so the segment $[u,z]$ lies in $X$.
For each $i$, the mean value theorem applied to $s\mapsto G_i(u+s(z-u))$ on $[0,1]$ gives a point $\zeta_i\in[u,z]$ with $G_i(z)-G_i(u)=\nabla G_i(\zeta_i)^{\top}(z-u)$.
Let $M$ be the matrix whose $i$th row is $\nabla G_i(\zeta_i)^{\top}$.
Its entries satisfy $M_{ij}=\partial_jG_i(\zeta_i)\in\mathbf{A}_{ij}$, because $G'(\zeta_i)\in\mathbf{A}$.
An interval matrix contains every real matrix whose entries lie in its entries, so $M\in\mathbf{A}$, and
\begin{equation}\label{eq:pts-mvf}
G(z)=G(u)+M(z-u),\qquad M\in\mathbf{A}.
\end{equation}

\emph{Step 2 (existence).}
Define $T(z)=z-CG(z)$ for $z\in X$.
By \eqref{eq:pts-mvf} with $u=y$, $T(z)=y-CG(y)+(I-CM)(z-y)$ with $M\in\mathbf{A}$.
Here $G(y)\in\mathbf{v}$ and $I-CM\in\mathbf{B}$ by (b), so $T(z)\in K\subseteq X$ by (c).
The map $T$ is continuous and maps the nonempty, compact, convex set $X$ into itself.
By Brouwer's fixed-point theorem there is $z^*\in X$ with $T(z^*)=z^*$, that is, $CG(z^*)=0$, and $z^*=T(z^*)\in K$.

\emph{Step 3 (nonsingularity).}
Choose $M\in\mathbf{B}$ with $|M_{ij}|=|\mathbf{B}_{ij}|$ for all $i,j$ (an endpoint of each entry), and fix $v\in\mathbf{v}$.
As $z$ ranges over $X$, its coordinates range independently over intervals of widths $w_j$.
So the $i$th component of $y-Cv+M(z-y)$ ranges over a closed interval of width $\sum_j|\mathbf{B}_{ij}|w_j$.
By (c), the closed interval $K_i$ contains this interval and lies in the open interval $\operatorname{int}X_i$ of length $w_i$.
Hence $\sum_j|\mathbf{B}_{ij}|w_j<w_i$ for every $i$, and
\[
q:=\max_i\frac{1}{w_i}\sum_j|\mathbf{B}_{ij}|w_j<1 .
\]
Let $D=\operatorname{diag}(w)$ and $A\in\mathbf{A}$.
By (b), $I-CA\in\mathbf{B}$, so $|(I-CA)_{ij}|\le|\mathbf{B}_{ij}|$ and
$\|D^{-1}(I-CA)D\|_\infty=\max_i w_i^{-1}\sum_j|(I-CA)_{ij}|w_j\le q<1$.
A matrix $I-P$ with $\|P\|_\infty<1$ is nonsingular (Neumann series), so $D^{-1}CAD=I-D^{-1}(I-CA)D$ is nonsingular, and so are $CA$, $C$ and $A$.

\emph{Step 4 (zero and uniqueness).}
$C$ is nonsingular and $CG(z^*)=0$, so $G(z^*)=0$.
If also $G(z)=0$ for some $z\in X$, then \eqref{eq:pts-mvf} with $u=z^*$ gives $0=M(z-z^*)$ with $M\in\mathbf{A}$.
$M$ is nonsingular by step 3, so $z=z^*$.

\emph{Step 5 (feasibility).}
At $x^*=(\xi,z^*)$ every row $i\in E$ has body value $t_i$, which is one of its bounds (its right-hand side if it is an equality), so the row holds; its expressions are defined by (a).
By (d), every other row and every bound holds at $x^*\in\{\xi\}\times X$, the objective and every expression are defined there, and every integer variable is fixed at an integer.
Hence $x^*\in\feas$.
Finally $z^*\in K$, so $f(x^*)$ lies in every enclosure of $f$ over $\{\xi\}\times K$.
\end{proof}

The implementations store a box $X$ with decimal endpoints, compute $\mathbf{A}$ by forward-mode interval differentiation of the OSIL expression trees, evaluate $\mathbf{v}$, $\mathbf{B}$ and $K=y-C\mathbf{v}+\mathbf{B}(X-y)$ in outward-rounded interval arithmetic, and test $K\subseteq\operatorname{int}X$.
The matrix $C$, a binary64 approximate inverse of a midpoint Jacobian, is used as an exact rational matrix and needs no error analysis.
If the enclosures are computed over an outward-rounded box $X'\supseteq X$ and the computed $K$ lies in $\operatorname{int}X$, the theorem applied to $X'$ places the unique zero of $X'$ in $K\subseteq X$, because $\operatorname{int}X\subseteq\operatorname{int}X'$.
The midpoint--radius test of the audit is a special case (\cref{lem:audit-krawczyk}, whose proof gives the reduction).
Condition (d) also shows that an equality row outside $E$ can pass only if it is constant on $\{\xi\}\times X$, for example because all its variables are fixed; no equality row passes by accident.

\subsection{Points in quadratic fields; proofs of the other results of \texorpdfstring{\cref{sec:points}}{Section 6}}\label{app:points-props}

Explicit exact points (\cref{sec:points-constructions}) rest on the following proposition.

\begin{proposition}[Points in $\Q$ or $\Q(\sqrt D)$]\label{prop:pts-exact}
Let $D>0$ be a rational number that is not the square of a rational.
\begin{itemize}
\item[(a)] Every element of $\Q(\sqrt D)$ has a unique form $p+q\sqrt D$ with $p,q\in\Q$, and its sign is decidable.
If $pq\ge0$, it is the sign of whichever of $p$ and $q$ is nonzero, and the element is zero if and only if $p=q=0$.
If $pq<0$, it is the sign of $p$ when $p^2>q^2D$ and the sign of $q$ when $p^2<q^2D$, and $p^2=q^2D$ cannot occur.
\item[(b)] Let $x$ be a point with coordinates in $\Q(\sqrt D)$, and let every row and bound of $\model$ be built from the coordinates and rational constants by $+$, $-$, $\times$, $\div$, integer powers and \texttt{sqrt}.
Evaluate every expression exactly in $\Q(\sqrt D)$; at a \texttt{sqrt} node with argument $a$, use a given candidate $s\in\Q(\sqrt D)$ after checking $s\ge0$ and $s^2=a$.
If every check passes, every denominator is nonzero, every row and bound comparison holds by the sign test in (a), and every integer variable has an integer value, then $x\in\feas$.
If the objective is built in the same way, $f(x)$ is the computed element of $\Q(\sqrt D)$.
\end{itemize}
The same holds row by row when each row involves coordinates from $\Q$ and at most one field $\Q(\sqrt{D_k})$.
\end{proposition}

\begin{proof}
(a) If $p+q\sqrt D=p'+q'\sqrt D$ with $q\ne q'$, then $\sqrt D=(p-p')/(q'-q)$ would be rational; so $q=q'$ and $p=p'$.
If $p,q\ge0$, then $p+q\sqrt D\ge0$, with equality if and only if $p=q=0$; the case $p,q\le0$ is symmetric.
If $pq<0$, then $p$ and $-q\sqrt D$ have the same sign, so $p-q\sqrt D\neq0$ has the sign of $p$.
From $(p+q\sqrt D)(p-q\sqrt D)=p^2-q^2D$, the sign of $p+q\sqrt D$ is the sign of $p$ times the sign of $p^2-q^2D$: the sign of $p$ if $p^2>q^2D$, and the sign of $-p$, which is the sign of $q$, if $p^2<q^2D$.
Since $q\ne0$, the equality $p^2=q^2D$ would make $D=(p/q)^2$ a rational square.

(b) $\Q(\sqrt D)$ is a field, so $+$, $-$, $\times$, integer powers (repeated products) and division by nonzero elements are exact operations on the pairs $(p,q)$, and part (a) decides whether a denominator is zero.
At a \texttt{sqrt} node with argument $a$, the value is the unique nonnegative real whose square is $a$; if $s\ge0$ and $s^2=a$, then $a\ge0$ satisfies the domain rule and the node equals $s$.
Every row and bound comparison is a sign test of a difference in $\Q(\sqrt D)$, decided by (a).
If all tests pass and the integer variables have integer values, every row, bound and integrality requirement holds and every expression is defined, so $x\in\feas$.
If the objective is built from the same operations, its value is the computed element.
For the row-by-row version apply the argument to each row in its own field.
\end{proof}

\begin{proof}[Proof of \cref{prop:pts-triangular}]
We show by induction on $t$ that $z^*_0,\dots,z^*_t$ are defined and lie in $Z_0,\dots,Z_t$.
For $t=0$ this holds because $Z_0=\{z_0\}$.
If it holds for $t-1$, then $(z^*_0,\dots,z^*_{t-1})\in Z_0\times\dots\times Z_{t-1}\subseteq\mathcal{D}_t$, so $z^*_t=\varphi_t(z^*_0,\dots,z^*_{t-1})$ is defined, and it lies in $Z_t$ because $\varphi_t(Z_0\times\dots\times Z_{t-1})\subseteq Z_t$.
The rows in $R_t$ involve only $z_0,\dots,z_t$, so by the defining property of $\varphi_t$ they hold exactly at $z^*$, with all their expressions defined.
Hence $z^*\in Z$ satisfies every row in $R_1\cup\dots\cup R_T$.
Every other row and bound holds at all points of $Z$, hence at $z^*$, the objective is defined there, and every integer variable equals its integral singleton value, so $z^*\in\feas$.
Finally $f(z^*)$ lies in every enclosure of $f$ over $Z$.
\end{proof}

The proof makes the symbolic step explicit: each row of $R_t$ must be rewritten correctly as a definition $z_t=\varphi_t(z_{<t})$.
Evaluating the defining rows over the boxes afterwards gives enclosures that contain $0$; this is consistent with exact satisfaction but does not prove it.
The safeguard is that the separately written implementations re-derived every defining row from the OSIL file and obtained the same definitions.

\begin{proof}[Proof of \cref{prop:pts-lindemann}]
Among the rows of \inst{lnts}$N$ are, for $i=0,\dots,N-1$,
\begin{gather*}
vx_{i+1}-vx_i-50h(\cos\theta_i+\cos\theta_{i+1})=0,\qquad
vy_{i+1}-vy_i-50h(\sin\theta_i+\sin\theta_{i+1})=0,\\
py_{i+1}-py_i-\tfrac h2(vy_i+vy_{i+1})=0,
\end{gather*}
with the fixed values $vx_0=vy_0=py_0=0$, $vx_N=45$ and $py_N=5$ (\cref{app:lnts}).
Suppose that a feasible point has algebraic $\theta_0,\dots,\theta_N$ and $h$.
Summing the $vx$ rows gives $45=100h\sum_{j=0}^{N}w_j\cos\theta_j$ with $w_0=w_N=\tfrac12$ and $w_j=1$ otherwise.
Hence $h\ne0$ and $\sum_jw_j\cos\theta_j=q:=9/(20h)$, which is algebraic.
Let $\mathcal{A}=\{|\theta_j|:\theta_j\ne0\}$, let $c_a=\sum_{j:|\theta_j|=a}w_j>0$ for $a\in\mathcal{A}$, and let $c_0=\sum_{j:\theta_j=0}w_j\ge0$.
Since $\cos$ is even and $\cos a=(e^{ia}+e^{-ia})/2$,
\[
(c_0-q)\,e^{0}+\sum_{a\in\mathcal{A}}\tfrac{c_a}{2}\,e^{ia}+\sum_{a\in\mathcal{A}}\tfrac{c_a}{2}\,e^{-ia}=0 .
\]
The exponents $0$ and $\pm ia$ ($a\in\mathcal{A}$) are distinct algebraic numbers, and the coefficients are algebraic.
By the Lindemann--Weierstrass theorem in the form of \citet[Chapter~1]{baker1975-transcendental-number-theory}, $e^{\beta_1},\dots,e^{\beta_k}$ are linearly independent over the algebraic numbers whenever $\beta_1,\dots,\beta_k$ are distinct algebraic numbers; hence every coefficient vanishes.
Since $c_a>0$, the set $\mathcal{A}$ is empty, so every $\theta_j=0$.
Then the $vy$ rows give $vy_i=vy_0=0$ for all $i$, and the $py$ rows give $py_N=py_0=0\ne5$, a contradiction.
\end{proof}

\subsection{Constructions by instance}\label{app:points-instances}

For each family we state the point, how its feasibility is proved, how its objective is enclosed and which implementations prove it, summarizing where \cref{app:families} gives the construction in full; ``separately written'' has the meaning of \cref{sec:semantics-protocol}.
Unless stated otherwise, $\Uprim$ is the upper end of the objective enclosure rounded up, as in \cref{tab:points-all}.

\begin{proposition}[Exactly feasible points for the 13 closures]\label{prop:pts-former}
For each instance of \cref{tab:points}, the point defined below exists and lies in $\feas$, and its objective lies in an enclosure of the width shown there.
For \inst{dtoc5}, \inst{lukvle10}, the \inst{chain} instances and the \inst{powerflow} instances, the upper end of this enclosure, rounded up, is the value $\Uprim$ of \cref{tab:closures}.
For the \inst{lnts} instances, $\Uprim$ is the attained optimum of \cref{thm:lnts-opt}, and the points of this proposition are a second, strictly suboptimal witness.
\end{proposition}

\begin{certbox}[Certificate for \cref{prop:pts-former}]
\certitem{Inputs}{OSIL files of the 13 instances (SHA-256 prefixes in \cref{tab:sem-hashes}). Point data: the 40-digit \inst{lnts} controls with their boxes; the \inst{dtoc5} decimal point; the 640-decimal \inst{lukvle10} seeds; the \inst{chain} generators (20-decimal $t_i$, $R_N$, root order); the \inst{powerflow} fixed values, square systems and 60-digit centres.}
\certitem{Computation}{Krawczyk tests (\cref{thm:pts-krawczyk}) on systems of 3 (\inst{lnts}) and 222, 262 and 270 (\inst{powerflow}) equations, with enclosures of all other rows and of the objective; exact evaluation of every row in $\Q$ (\inst{dtoc5}) or $\Q(\sqrt{R_N})$ (\inst{chain}); a forward recursion with enclosures (\inst{lukvle10}). Arithmetic: \arith{I} (mpmath at 110, 750 and 80 digits) in the first implementation for \inst{lnts}, \inst{lukvle10} and \inst{powerflow}; \arith{E} for \inst{dtoc5} and \inst{chain} and in every second implementation.}
\certitem{Data reading}{exact decimals in every reader.}
\certitem{Implementations}{each point is proved by the first implementation and again by a separately written implementation with its own OSIL reader that uses only integers and rationals with explicit series remainders: dyadic intervals with alternating-series bounds for $\sin$ and $\cos$ (\inst{lnts}); \texttt{Fraction} arithmetic, plus a third exact reader (\inst{dtoc5}); integer fixed point with integer logarithm and exponential (\inst{lukvle10}); its own arithmetic in $\Q(\sqrt D)$ (\inst{chain}); dyadic intervals at $2^{-320}$ with Taylor bounds for $\sin$ and $\cos$ (\inst{powerflow}). Either implementation alone proves the claim.}
\certitem{Evidence level}{\evid{hand} (\cref{thm:pts-krawczyk,prop:pts-exact,prop:pts-triangular}) $+$ \evid{stored}.}
\certitem{Replay}{Tier~1; seconds to minutes per family (\cref{app:repro}).}
\end{certbox}

\paragraph{\inst{lnts50}--\inst{lnts400}.}
Each instance has two exactly feasible points.
The first, which gives $\Uprim$, is the attaining point of \cref{thm:lnts-opt}; its existence rests on the intermediate value theorem for a scalar root, and two integer-only implementations enclose its objective in exact rational arithmetic \arith{E}, with width below $\sci{1.11}{-91}$ (\cref{app:lnts}).
The second is the point of \cref{prop:pts-former}, a second, strictly suboptimal witness (\cref{rem:lnts-stored}).
Its controls $\theta_1,\dots,\theta_{N-1}$ are fixed at 40-digit decimal rationals that round the linear-tangent law.
Given $(\theta_0,\theta_N,h)$, the rows define $vx$, $vy$, $px$ and $py$ by $v_{i+1}=v_i+50h(\phi(\theta_i)+\phi(\theta_{i+1}))$ and $p_{i+1}=p_i+\tfrac h2(v_i+v_{i+1})$, with $(v,p,\phi)=(vx,px,\cos)$ or $(vy,py,\sin)$, starting from $0$.
Every row has coefficient $1$ on its next state, so all $4N$ rows hold identically (\cref{prop:pts-triangular}).
The fixed values that remain give $G(\theta_0,\theta_N,h)=(vx_N-45,\ vy_N,\ py_N-5)=0$; $px_N$ is free.
\Cref{thm:pts-krawczyk} with $m=3$ is applied on a box around a 75-digit centre with radius $10^{-50}$ in $\theta_0$ and $\theta_N$ and $10^{-52}$ in $h$, with the Jacobian obtained by forward-mode differentiation of the recursion and $C$ the inverse of the midpoint Jacobian.
On the box $|\theta_j|\le0.9523$ and $h>0$, so all bounds hold.
The first implementation uses mpmath intervals at 110 digits; the separately written one re-derived the elimination from the OSIL file and repeated the test in mpmath at 120 digits and in integer dyadic arithmetic.

\paragraph{\inst{dtoc5}.}
The model has the rows $-\tfrac{1}{50000}u_t+y_t-y_{t+1}+\tfrac{1}{12500}y_t^2=0$ (decimal strings \texttt{2e-5} and \texttt{8e-5}) for $t=0,\dots,49998$, the fixed value $y_0=1$ and no other bounds (\cref{app:dtoc5}).
The states $y_1,\dots,y_{49999}$ of a Newton solution polished at 50 digits are rounded to 30 decimals, and each control is \emph{defined} by its row, $u_t:=50000\,(y_t-y_{t+1})+4y_t^2$, a decimal with at most 60 places.
So every row holds identically, and the objective is an exact rational whose value gives $\Uprim$.
Exact checks with separate readers confirm every row, the fixed value $y_0=1$ and the objective: the first code's generic checker with a cross-check by another reader, a separately written implementation with its own reader, and a third exact reader.
As a negative control, adding $10^{-60}$ to one control makes exactly its row fail.

\paragraph{\inst{lukvle10}.}

\begin{remark}[Seeds for \inst{lukvle10}]\label{rem:lukvle10-seeds}
Every row of \inst{lukvle10} reads $-x_j+3x_{j+1}-2x_{j+2}-2x_{j+1}^2=-1$ ($j=0,\dots,997$), and all 1000 variables are free.
Row $j$ defines $x_{j+2}=(1-x_j+3x_{j+1}-2x_{j+1}^2)/2$, and the objective is defined at every real point, because its \texttt{power} nodes have exponents of the form $x_k^2+1>0$ (\cref{def:sem-model}).
Any two rational seeds $x_0,x_1$ therefore define a rational, exactly feasible point (\cref{prop:pts-triangular}).
The denominators roughly square at every step, so the point cannot usefully be printed.
Useful seeds must also be very precise.
At the fixed point $-1/\sqrt2$ of the recursion, its linearization has the eigenvalues $2.731$ and $0.183$, so seed errors grow by a factor of about $2.73^{998}\approx10^{435}$ along the chain.
In a floating-point experiment at 1500 digits, the 15-digit seeds of MINLPLib's listed point p5 drive $|x_j|$ above $10$ at index $40$.
Seeds that agree with a numerical local minimizer to 50--420 decimals give objective values between $352.89$ and $353.01$; seeds with 440 decimals give a value within $\sci{2}{-14}$ of the local minimizer's, and 460 or more decimals reproduce it to 20 digits.
Our seeds are the local minimizer's values rounded to 640 decimals; the seeds and the recursion define our point.
MINLPLib's listed point p5, whose rows hold to $\sci{3.48}{-15}$, evaluates about $\sci{5.1}{-13}$ above the objective of our point.
\end{remark}

The first implementation encloses all coordinates in mpmath intervals at 750 digits; the smallest $|x_k|$ is $0.26003$, so every \texttt{power} base $x_k^2$ is at least $0.0676$, and the objective is enclosed through $\exp(b\ln a)$ with positive bases $a$.
The objective enclosure gives $\Uprim=352.2380254064957$.
A further implementation on the first side uses integer fixed point with an $\operatorname{atanh}$-series logarithm and a Taylor exponential; the separately written one uses integer fixed point at $2^{-3400}$ with its own integer logarithm and exponential.
Both reproduce the enclosure, so the primal side of \inst{lukvle10} does not depend on mpmath.
The point agrees with the numerical local minimizer to about 205 digits.
This is evidence of local stationarity only; global optimality of the point is not proved.

\paragraph{\inst{chain50}--\inst{chain400}.}
The model has variables $x_0,\dots,x_N$ (heights) and $u_0,\dots,u_N$ (slopes), all free except $x_0=1$ and $x_N=3$, the rows $x_{i+1}-x_i-\eta(u_i+u_{i+1})=0$ for $i<N$, and the length row $\eta\sum_{i<N}(s_i+s_{i+1})=4$, where $\eta=1/(2N)$ and $s_i=\sqrt{1+u_i^2}$ are \texttt{sqrt} nodes.
Let $\omega_0=\omega_N=1$ and $\omega_i=2$ otherwise.

\begin{lemma}[Rational parametrization of \inst{chain}]\label{lem:pts-chain}
For $t\in\R^{N+1}$ with $t>0$, put $u_i=(t_i-t_i^{-1})/2$, let $x_0=1$ and define $x_1,\dots,x_N$ by the linear rows.
Then $s_i=(t_i+t_i^{-1})/2$ for every $i$, and $(x,u)$ is exactly feasible if and only if
\[
\sum_{i=0}^N\omega_it_i=12N\qquad\text{and}\qquad\sum_{i=0}^N\omega_i/t_i=4N .
\]
Every exactly feasible point arises in this way, from $t_i=u_i+\sqrt{1+u_i^2}$.
\end{lemma}

\begin{proof}
Put $\sigma_i=(t_i+t_i^{-1})/2$.
Then $\sigma_i>0$ and $\sigma_i^2-u_i^2=1$, so $\sigma_i=\sqrt{1+u_i^2}=s_i$.
Summing the linear rows gives $x_N=1+\eta\sum_{i<N}(u_i+u_{i+1})=1+\eta\sum_i\omega_iu_i$, so $x_N=3$ if and only if $\sum_i\omega_iu_i=4N$; the length row reads $\sum_i\omega_is_i=8N$.
Since $t_i=s_i+u_i$ and $t_i^{-1}=s_i-u_i$, these two conditions are equivalent to $\sum_i\omega_it_i=12N$ and $\sum_i\omega_i/t_i=4N$.
All other rows hold by the definition of $x$, and every expression is defined.
Conversely, for any real $u_i$, $t_i=u_i+\sqrt{1+u_i^2}>0$ satisfies $u_i=(t_i-t_i^{-1})/2$, and the linear rows determine $x$ from $x_0=1$.
\end{proof}

Starting from a floating-point local minimizer, we fix $t_i$ for $i\notin\{1,N-1\}$ at 20-decimal rationals.
With $\alpha=(12N-\sum'\omega_it_i)/2$ and $\beta=(4N-\sum'\omega_i/t_i)/2$, where $\sum'$ omits $i=1$ and $i=N-1$, the lemma requires $t_1+t_{N-1}=\alpha$ and $t_1^{-1}+t_{N-1}^{-1}=\beta$, that is, $t_1t_{N-1}=\alpha/\beta$.
So $t_1$ and $t_{N-1}$ are the roots of $T^2-\alpha T+\alpha/\beta$; we take $t_1$ to be the smaller root.
Both roots are positive because $\alpha>0$, $\alpha/\beta>0$ and the discriminant $\alpha^2-4\alpha/\beta$ is positive, as the checks verify.
Writing the discriminant as $n/d$ in lowest terms and $R_N=nd$, which is not a square, every coordinate lies in $\Q(\sqrt{R_N})$; $R_N$ has 1,928, 3,755, 7,323 and 14,507 digits.
Feasibility is decided by \cref{prop:pts-exact}: every row and both fixed values hold exactly in $\Q(\sqrt{R_N})$, and each \texttt{sqrt} node is accepted because the candidate $(t_i+t_i^{-1})/2$ is nonnegative and squares to $1+u_i^2$.
The objective is enclosed to 40 decimals with integer square roots.
The first code's build and verification scripts are separate programs, and the separately written implementation repeated every check.
SCIP's floating-point solution check accepts binary64 roundings of the points at feasibility tolerance $10^{-13}$ (evidence only).
For \inst{chain50}, the other root order also gives an exactly feasible point, with a much larger objective ($6.326$), so the root order is part of the definition.
We do not know whether a rational feasible point exists near the optimum.

\paragraph{\inst{powerflow0030p}, \inst{powerflow0039p}, \inst{powerflow0039r}.}
\Cref{app:powerflow-points} gives the construction and the two implementations.
Starting from MINLPLib's p1, rows with slack below $10^{-9}$ are treated as active; the active slacks are at most $\sci{6.8}{-13}$ and the next slacks at least $\sci{4.3}{-4}$.
Compared with the scheme of \citet{kearfott1998-on-proving-existence-of-feasible}, the fixed coordinates are exact rationals, active multi-variable inequality rows are imposed as equalities, and the box is far smaller (radius $10^{-45}$).
The free variables of the exact points differ from those of p1 by at most $\sci{1.4}{-14}$, $\sci{2.5}{-13}$ and $\sci{7.8}{-12}$, and their objectives exceed those of p1 by $\sci{2.87}{-12}$, $\sci{4.64}{-12}$ and $\sci{1.83}{-10}$.
SCIP's floating-point check accepts the centres rounded to binary64 at feasibility tolerance $10^{-9}$ (evidence only).

\paragraph{\inst{optcdeg2}.}
\Cref{prop:optcdeg2-point} gives the point and its proof (\cref{app:optcdeg2-point}): the controls are $\pm1/5$ on the three arcs, one switching control is a stored binary64 value, and the other is bracketed by the intermediate value theorem applied to the terminal condition $v_N=0$, in integer interval arithmetic and again, separately written, in mpmath.
The point concerns MINLPLib's coefficient $0.2$ in the velocity rows, not CUTEst's $0.05$ (\cref{app:optcdeg2-model}).

\paragraph{\inst{catmix100}--\inst{catmix800}.}
With $a=1/(2N)$ and the stored coefficient $c$, the rows of \eqref{eq:catmix-model} are $P(u_{i+1})x_{i+1}=Q(u_i)x_i$ for $i<N$, with $x_0=(1,0)$, free states and controls $u\in[0,1]^{N+1}$.
Since $\det P(u)\ge1+a$ for $u\in[0,1]$ (\cref{lem:catmix-nonneg}), any controls in $[0,1]$ define the states uniquely by $2\times2$ blocks, and every row holds exactly (\cref{prop:pts-triangular,lem:catmix-reduction}).
For binary64 controls the states and the objective $x_{1,N}+x_{2,N}-1$ are exact rationals, and $\Uprim$ is the exact objective rounded up.
We use the controls of the points computed by the first code for $N=100,200,400$ and those of the policy of the dual certificate for $N=800$.
For each $N$, at least two separately written implementations with their own readers computed the objective in exact rational arithmetic; one of them solves each pair of rows as a $2\times2$ linear system without using the structure of $P$ and $Q$.
For $N=400$ the first code's 60-digit mpmath interval simulation agrees.

\paragraph{\inst{camshape100}--\inst{camshape800}.}
The point is the envelope $E$ of \cref{lem:camshape-envelope} with the slope variables $d_i=E_{i+1}-E_i$; all coordinates are rationals computed by exact recurrences.
\Cref{lem:camshape-envelope}(b) proves feasibility, and two separately written exact implementations confirm it by evaluating every OSIL row and bound in rational arithmetic.
The objective equals $v_n$ exactly, so $\Uprim$ is the ceiling display of the exact optimum.
A binary64 rounding of $E$ violates rows by $\sci{2.8}{-16}$ to $\sci{4.8}{-16}$; the exact point is the witness.

\paragraph{The six small instances.}
\Cref{app:small} gives each construction and its proof.
The \inst{ex6_2_5} and \inst{ex6_2_7} points are rational, so their rows and bounds hold exactly; their objectives are enclosed in mpmath intervals and, without mpmath, in rational intervals.
The \inst{pricing050} point is a saved 17-digit point whose five rows hold, by mpmath interval evaluation in two separately written implementations with their own readers, with slack at least $\sci{3.49}{-15}$; its bounds hold exactly and its linear objective is exact.
The \inst{etamac} point is defined forward from saved 17-digit decisions (\cref{prop:pts-triangular}); mpmath intervals prove its bounds and enclose its objective.
The \inst{pindyck} point fixes the prices; each supply row $s-\hat a-\hat b\,e^{-Ks}=0$ with $\hat b>0$ and $K>0$ has a strictly increasing left side with limits $\mp\infty$, so it defines the supply uniquely, and a one-dimensional interval test in mpmath encloses it in two implementations (\cref{lem:pindyck-reduction}).
The \inst{hvycrash} point fixes $c_k=0.08$ and is defined backwards from $\theta_{50}=3$: the rows of stage $k$ give $r_k=(-\cos\theta_k/D(0.08))^{1/2}$, then $\theta_{k-1}$ and $s_k=-kh$, which are defined because every $\cos\theta_k<-0.85$; \cref{prop:hvycrash-identity} proves its feasibility in a written proof, two codes checked this witness in mpmath intervals and a third code a second witness.
The stored decimal vectors of the \inst{hvycrash}, \inst{etamac} and \inst{pindyck} points are not exactly feasible; the points are the real points defined by the recursions.

\paragraph{\inst{eg_int_s}, \inst{eg_disc_s}, \inst{eg_disc2_s}.}
\Cref{app:eg-primal} lists the points, their slacks and the tests of the exponential below.
All 28 rows are inequalities.
The objective variable, whose value is the objective, is a 20-digit decimal rounded up from a high-precision evaluation of the largest of the 24 right-hand sides; this rounding gives the smallest margins, $\sci{5.57}{-20}$, $\sci{7.80}{-20}$ and $\sci{7.49}{-20}$.
A library-free implementation with its own OSIL reader proves all 28 rows in dyadic intervals (integers times $2^{-256}$, with floor and ceiling rounding), bounding $\exp$ by the alternating Taylor series of $\exp(-y)$ for $0\le y\le2^{-8}$ followed by outward squaring.
A separate agent session reran this check, read it in full and confirmed the three points; its tests support, and do not replace, the alternating-series proof.
Two separately written mpmath interval implementations with their own readers give the same margins, and six perturbed copies of the \inst{eg_int_s} point are rejected.
MINLPLib's listed points p1 are not claimed to be feasible.

\paragraph{\inst{waterno2_06}--\inst{waterno2_24}.}
The construction starts from MINLPLib's p4 for \inst{waterno2_06} and from our own window-reoptimized points for the other four instances.
First, the binaries are fixed, the variables forced by single rows are pinned (an off pump has flow $0$ and minimal speed), and the big-M rows of running pumps become equalities.
Second, the flows of the running pumps of one station are aliased, because a shared speed and head force equal flows.
Third, the linear rows are solved exactly, with active level bounds and the horizon row imposed by a minimum-norm exact rational correction of the station flows; the largest correction per instance lies between $\sci{3.0}{-16}$ and $\sci{8.8}{-9}$.
Fourth, each running station's pump-curve row becomes a quadratic $Aw^2+Bw+C=0$ in its speed $w$ with rational coefficients; the speed is the root nearest the numerical speed, isolated in a rational interval of width $\sci{2}{-45}$ by a sign change, and written $w=(-B+\sigma\sqrt{B^2-4AC})/(2A)$ with the sign $\sigma$ fixed by that interval.
The points have 9, 18, 28, 48 and 63 such irrational speeds.
Finally, the cost variables are set to rational upper bounds of the power expressions on a $10^{-30}$ grid, so the objective is an exact rational.
Every coordinate lies in $\Q$ or in one quadratic field, and no row mixes two fields, so \cref{prop:pts-exact} decides every row, bound and integrality requirement exactly.
Three exact checkers ran: two of the first code, one over the quadratic fields and one that decides signs by comparing $p^2$ with $q^2D$, and a separately written evaluator over the multi-quadratic field $\Q(\sqrt{D_1},\dots,\sqrt{D_k})$ with its own reader, whose mutation tests reject perturbations of $10^{-30}$.
The \inst{waterno2_06} point lies within $\sci{3.1}{-13}$ of p4; it is MINLPLib's point made exact, not our own.
The point p4 itself violates row \texttt{e94} by $\sci{1.08}{-11}$.
The binary64 reading of these points is treated in \cref{app:waterno2-model}.

\paragraph{\inst{ann_cumene_tanh}.}
The five inputs are fixed at the rationals given by the 16-digit decimals of a computed point.
Every other variable is defined by one of the 789 equality rows, in an acyclic order, in which it is the only unknown and appears linearly with a nonzero coefficient: a constant, or, in rows \texttt{e749} and \texttt{e782}--\texttt{e787}, a previously defined variable ($x_{753}$, $x_{746}$ or $x_{766}$) whose enclosure excludes $0$ (\cref{prop:pts-triangular}).
The inequality row \texttt{e789} holds with margin $\sci{9.979}{-13}$, and the bounds of all 728 bounded variables hold.
The tightest are $x_{647}\ge-1$, with margin $\sci{1.25}{-12}$, and the saturated neurons $x_{650}$ and $x_{657}$, whose enclosures touch $\pm1$; this still proves their non-strict bounds.
The first implementation uses mpmath intervals at 60 digits; a separately written implementation with its own reader uses rational intervals with a Taylor exponential and Lagrange remainder, and its enclosures (width $\sci{2.8}{-89}$) lie inside those of the first.
The 30- and 40-digit decimal roundings of the point are not exactly feasible.

\paragraph{KAN instances.}
The points of the six KAN instances of our set lie in $\RP$, the OSIL model without the partition-of-unity rows (\cref{app:annkan-kan}), because the stored models have no exactly feasible point (\cref{prop:kan-infeasible}).
The inputs are fixed at the binary64 values of computed minimizers.
For each edge, the binary variable of the first knot interval whose big-M rows hold over the whole enclosure of the edge argument is set to $1$, so the one-hot rows hold exactly.
Every other equality row except the partition-of-unity rows defines one variable as in \cref{prop:pts-triangular}.
The big-M rows hold exactly in the first layer, where the arguments are rational, and by enclosures in the second, with smallest margins between $\sci{7.3}{-5}$ and $\sci{6.5}{-2}$ across the six instances; every OSIL variable bound holds, including the bounds that $\Rnet$ drops.
The first implementation uses mpmath intervals at 60 digits; a separately written implementation uses rational intervals and encloses each objective with width at most $\sci{3.2}{-89}$.
The values $\Uprim$ are those of \cref{tab:kan}.

\subsection{All 43 points}\label{app:points-table}

\Cref{tab:points-all}, generated from the number file of the archive (\cref{app:repro}), gives for each family the construction of its points and the arithmetic of their proofs.

% Generated by paper-open-minlplib/data/make_points_table.py from data/numbers.json.
% Do not edit by hand; rerun the script. Requires booktabs and array;
% \inst{} takes instance names with unescaped underscores (macros.tex).
\begin{table}[htbp]
\centering
\scriptsize
\setlength{\tabcolsep}{3pt}
\caption{The exactly feasible point of each instance family; the values $U$ are those of \cref{tab:closures,tab:unclosed,tab:kan}.
Construction (\cref{sec:points-constructions}): \emph{exact point}, explicit exact point; \emph{triangular}, triangular definitions; \emph{existence proof}, interval existence proof; \emph{interior point}, strictly interior point.
Arithmetic: tags of the first and of the separately written second proof (\arith{E} exact rational or algebraic, including integer and dyadic interval arithmetic; \arith{I} mpmath intervals).
No mpmath: a computer proof without mpmath exists.
\textsuperscript{a}Points of $\RP$; the stored models have no exactly feasible point (\cref{prop:kan-infeasible}).
\textsuperscript{b}The existence proof is a written proof (\cref{prop:hvycrash-identity}; evidence level \evid{hand}); the computer checks of the witnesses use mpmath.
\textsuperscript{c}Maximization: $U$ is the lower end, rounded down.}
\label{tab:points-all}
\begin{tabular}{@{}>{\raggedright\arraybackslash}p{3.0cm}>{\raggedright\arraybackslash}p{1.7cm}>{\raggedright\arraybackslash}p{5.2cm}>{\raggedright\arraybackslash}p{3.2cm}>{\raggedright\arraybackslash}p{1.4cm}@{}}
\toprule
instances & construction & point & arithmetic & no mpmath \\
\midrule
\inst{lnts50}--\inst{lnts400} & existence proof & attaining point of \cref{thm:lnts-opt} (scalar root); a second point by a Krawczyk test (\cref{tab:points}) & \arith{E}; \arith{E} & yes \\
\inst{dtoc5} & exact point & rational point: states rounded to 30 decimals, controls from the rows & \arith{E} (three readers) & yes \\
\inst{optcdeg2} & existence proof & controls $\pm1/5$ with one switching control fixed and one bracketed (intermediate value theorem) & \arith{E}; \arith{I} & yes \\
\inst{lukvle10} & triangular & seeds with 640 decimals, forward recursion (\cref{rem:lukvle10-seeds}) & \arith{I}; \arith{E} & yes \\
\inst{chain50}--\inst{chain400} & exact point & point in $\Q(\sqrt{R_N})$ (\cref{lem:pts-chain}) & \arith{E}; \arith{E} & yes \\
\inst{catmix100}--\inst{catmix400} & triangular & binary64 controls of a computed run, exact rational states & \arith{E}; \arith{E} & yes \\
\inst{catmix800} & triangular & binary64 controls of the dual certificate\textquoteright s policy, exact rational states & \arith{E}; \arith{E} & yes \\
\inst{camshape100}--\inst{camshape800} & exact point & rational envelope point (\cref{lem:camshape-envelope}) & \arith{E}; \arith{E} & yes \\
\inst{ex6_2_5}, \inst{ex6_2_7} & exact point & rational point; linear rows exact; objective enclosed & rows \arith{E}, objective \arith{I}; \arith{E} & yes \\
\inst{pricing050}\textsuperscript{c} & interior point & rational interior point (17 digits) & \arith{I} (two codes) & no \\
\inst{etamac} & triangular & 17-digit decisions, forward definitions & \arith{I} (one code) & no \\
\inst{pindyck} & existence proof & 17-digit prices; states by recursion, one-dimensional interval test per supply row & \arith{I} (two codes) & no \\
the three \inst{powerflow} instances & existence proof & Krawczyk test on an active-set square system near p1 & \arith{I}; \arith{E} & yes \\
\inst{hvycrash} & triangular & backward definitions from $\theta_{50}=3$ (\cref{prop:hvycrash-identity}) & written proof; \arith{I} (two codes for the first witness, a third for a second) & no\textsuperscript{b} \\
the three \inst{eg} instances & interior point & decimal point, objective variable rounded up; strictly interior rows & \arith{E}; \arith{I} & yes \\
\inst{waterno2_06} & exact point & MINLPLib\textquoteright s p4 made exact; rational or one quadratic field per row & \arith{E} (three checkers) & yes \\
\inst{waterno2_09}--\inst{waterno2_24} & exact point & algebraic point: rational or one quadratic field per row & \arith{E} (three checkers) & yes \\
\inst{ann_cumene_tanh} & triangular & rational inputs, forward definitions & \arith{I}; \arith{E} & yes \\
the six KAN instances\textsuperscript{a} & triangular & point of $\RP$: binary64 inputs, forward definitions & \arith{I}; \arith{E} & yes \\
\bottomrule
\end{tabular}
\end{table}

